\documentclass[11pt]{article}

\usepackage{../ppackage}
\usepackage{bbm}
\usepackage{import}
\usepackage{refstyle}






\usepackage{tikz}
\usetikzlibrary{arrows.meta,patterns}
\usetikzlibrary{ipe} % ipe compatibility library

\usepackage{../../../Notes/tikzit}
\usepackage{../../../Notes/utf8math}

\input{../../../Notes/TIKZ/digraph.tikzstyles}



\usepackage{url}


\newcommand{\solution}{
\bigskip\noindent
	\textbf{Solution: \\}}
	


\DeclareMathOperator{\size}{size}
\DeclareMathOperator{\conv}{conv}
\newcommand{\SV}{\mathrm{SV}}
\newcommand{\bigO}{O}
\newcommand{\cut}{\mathrm{cut}}
\newcommand{\LLL}{\mathrm{LLL}}
\newcommand{\setR}{\mathbb{R}}
\newcommand{\setZ}{\mathbb{Z}}
\newcommand{\setQ}{\mathbb{Q}}
\newcommand{\setC}{\mathbb{C}}
\newcommand{\setN}{\mathbb{N}}
\newcommand{\wt}[1]{\widetilde{#1}}
\newcommand{\opt}{{\sc 0/1-opt}\xspace}
\newcommand{\aug}{{\sc 0/1-aug}\xspace}
\newcommand{\psep}{{\sc 0/1-psep}\xspace}
\newcommand{\sep}{{\sc 0/1-sep}\xspace}
\newcommand{\fopt}{{\sc 0/1-testopt\xspace} }

\newcommand{\hpp}{\mathrm{HPP}}
\newcommand{\nodes}{\mathcal{V}}
\newcommand{\vol}{\mathrm{vol}}
\newcommand{\diag}{\mathrm{diag}}
\newcommand{\arcs}{\mathcal{A}}
\newcommand{\edges}{\mathcal{E}}
\newcommand{\paths}{\mathscr{P}}
\newcommand{\cycles}{\mathcal{C}}




\newcommand{\K}{{\mathcal K}}
\newcommand{\A}{{A}}
\newcommand{\B}{{B}}
\newcommand{\T}{\mathscr{T}}
\newcommand{\eE}{\mathscr{E}}
\newcommand{\eS}{\mathscr{S}}
\newcommand{\eP}{\mathscr{P}}
\newcommand{\eM}{\mathscr{M}}



\newcommand{\transp}{^{\mathrm{T}}}

\newcommand{\smallmat}[1]{\left( \begin{smallmatrix} #1 \end{smallmatrix}\right)}

\newcommand{\mat}[1]{ \begin{pmatrix} #1 \end{pmatrix}}
\newcommand{\smat}[1]{ \big(\begin{smallmatrix} #1 \end{smallmatrix}\big)}

\newcommand{\pc}{\mathscr{P}}
\newcommand{\ob}{\mathscr{O}}
\newcommand{\odds}{\mathscr{W}}
\newcommand{\up}{\mathscr{U}}
\newcommand{\ef}{\mathscr{F}}
\newcommand{\eh}{\mathscr{H}}
\newcommand{\ev}{\mathscr{V}}
\newcommand{\ec}{\mathscr{C}}
\newcommand{\eu}{\mathscr{U}}

\newcommand{\lex}{\mathrm{lex}}

\renewcommand{\leq}{\leqslant}
\renewcommand{\geq}{\geqslant}









\newcommand{\linhull}{\mathrm{lin.hull}}
\newcommand{\affhull}{\mathrm{affine.hull}}
\newcommand{\charcone}{\mathrm{char.cone}}
\newcommand{\cone}{\mathrm{cone}}
\newcommand{\rank}{\mathrm{rank}}
\newcommand{\wb}[1]{\overline{#1}}



\usepackage{enumerate}

      
\institute{\'Ecole Polytechnique F\'ed\'erale de Lausanne}
\lecture{Discrete Optimization}
\faculty{Prof. Friedrich Eisenbrand}
\term{Spring 2026}
\publishdate{May 19, 2026}
\duedate{ }
\problemset{Problem Set~12 (Solutions)}

\begin{document}
\makeheader

\begin{enumerate}[1)]

\item \label{Ex:Node-Pot} Let $G = (V,A)$ be a directed graph and $c: A \to \mathbb{R}$ be a length function. Consider the following set of linear inequalities with variables $\pi$:
  \begin{equation}\label{eq:pot5}
    \forall (u,v) \in A: \pi(v) \leq \pi(u) + c_{uv}.
  \end{equation}
  \begin{enumerate}[i)]
  \item   Show that this system of inequalities has a feasible solution if and only if $G$ does not possess a cycle of negative weight.
  \item  Interpret a cycle of negative weight as a Farkas certificate of infeasibility.
  \item Let $s \in V$ be a vertex from where each other vertex of the graph can be reached.  Consider the linear program which maximizes $\sum_{v \in V} \pi(v) $ subject to \ref{eq:pot5} and the additional constraint $\pi(s) = 0$. Show that the linear program has a unique optimal solution $\pi^*$.
  
  \emph{Hint: What is $\pi^*$ in terms of distances from $s$? }
\end{enumerate}

\begin{solution}
\begin{enumerate}[i)]
\item We first prove $\Rightarrow$. Suppose that there is a feasible solution $\pi: V \to \mathbb{R}$. For any cycle $C=v_0\rightarrow v_1\rightarrow \cdots \rightarrow v_{k-1}\rightarrow v_0$ in $G$, 
apply equation (\ref{eq:pot5}) to each arc of $C$ to get
\[
\pi(v_{i+1}) \leq \pi(v_i) + c_{v_i v_{i+1}} \qquad i=0,\ldots,k-1,
\]
where indices are modulo $k$. Summing up all to get
\[
\sum_{i=0}^{k-1} \pi(v_{i+1})
\leq
\sum_{i=0}^{k-1} \pi(v_i) + \sum_{i=0}^{k-1} c_{v_i v_{i+1}} \quad\Rightarrow\quad 0 \leq \sum_{i=0}^{k-1} c_{v_i v_{i+1}}.
\]

Next we prove $\Leftarrow$. Assume that $G$ has no negative-weight directed cycle. Add a new vertex $r$ and arcs $(r,v)$ of length $0$ to every $v \in V$. Let $d(v)$ be the minimum distance from $r$ to $v$ in this new graph. Since there is no negative cycle, all these distances are well-defined. For every arc $(u,v) \in A$, by definition of $d$, we have $d(v) \leq d(u) + c_{uv}$. Hence $d(v)$ is a feasible solution.

\item Let $C$ be a cycle of negative weight. 
Define $\lambda\in \mathbb{R}^{A}_{\geq 0}$ such that $\lambda_{uv}=1$ if $(u,v)\in C$ and $\lambda_{uv}=0$ otherwise. Then
$$\sum_{(u,v)\in A} \lambda_{uv} (\pi(v)-\pi(u)) = \sum_{(u,v)\in C} (\pi(v)-\pi(u)) = 0,$$ 
and 
$$\sum_{(u,v)\in A} \lambda_{uv} c_{uv} = \sum_{(u,v)\in C} c_{uv} < 0.$$
Thus $\lambda$ is a Farkas certificate of infeasibility.

\item For each vertex $v \in V$, let $\pi^*(v)$ be the distance of the shortest path from $s$ to $v$, which is finite since every vertex $v$ is reachable from $s$. Since there are no negative cycles, the distances are well-defined; also they satisfy
\[
\pi^*(v) \leq \pi^*(u)+c_{uv} \qquad \forall \;(u,v)\in A,
\]
and $\pi^*(s)=0$. Hence $\pi^*(v)$ is feasible for the linear program.

Now let $\pi: V \to \mathbb{R}$ be any feasible solution which satisfies \ref{eq:pot5} with $\pi(s)=0$. For any vertex $v \in V$, consider a directed $s$--$v$ path
$s=v_0 \to v_1 \to \cdots \to v_k=v$, one has
\[
\pi(v) \leq \pi(s)+\sum_{i=0}^{k-1} c_{v_i v_{i+1}}
= \sum_{i=0}^{k-1} c_{v_i v_{i+1}},
\]
which implies that $\pi(v) \leq \pi^*(v)$. Thus $\pi^*$ maximizes $\sum_{v \in V} \pi(v)$. 
Eventually we prove the uniqueness. Suppose that there exists another optimal solution $\pi'\neq \pi^*$. Then for some $v \in V$, we have $\pi'(v) < \pi^*(v)$. However, this implies that $\sum_{v \in V} \pi'(v) < \sum_{v \in V} \pi^*(v)$, contradicting the optimality of $\pi'$. 
\end{enumerate}
\end{solution}

\item \label{ex:count:paths} Let $G = (V,A)$ be a directed graph. The \emph{directed adjacency matrix} $A \in \{0,1\}^{V \times V}$ is  defined as 
\[
a_{uv} =
\begin{cases}
1, & \text{if there is an arc } u\to v,\\
0, & \text{otherwise.}
\end{cases}
\]
So \(A\) encodes exactly which directed edges, or arcs, are present in the graph.
\begin{enumerate}[i)] 
\item Show by induction that for  $k \in \mathbb{N}_+$ that  $A^k$ encodes the number of directed walks of length $k$ from $u$ to $v$. In other words
  \[
a^k_{uv} =  \text{ number of directed walks of length $k$ from $u$ to $v$.} 
\]
\item Show that $A^k$ can be computed with $O(\log k \cdot n^{2.805})$ arithmetic operations.
\end{enumerate}

\begin{solution}
\begin{enumerate}[i)]
\item For $k=1$ the claim is trivial. Suppose that the claim holds for some $k>1$. We will show that it holds for $k+1$. One has
\[
(A^{k+1})_{uv}=(A^kA)_{uv}=\sum_{w\in V}(A^k)_{uw}a_{wv}.
\]
By the induction hypothesis, $(A^k)_{uw}$ is the number of directed walks of length $k$ from $u$ to $w$. Multiplying by $a_{wv}$ keeps exactly those walks that can be extended by the arc $w\to v$. Hence the summand counts the directed walks of length $k+1$ from $u$ to $v$ whose penultimate vertex is $w$. Summing over all $w\in V$ counts all directed walks of length $k+1$ from $u$ to $v$. This proves the induction step.

\item Compute $A^k$ by repeated squaring. Writing $k$ in binary representation, one needs $O(\log k)$ matrix multiplications.
Using the algorithm for multiplying two $n\times n$ matrices with arithmetic complexity $O(n^{2.805})$. Thus the total number of arithmetic operations is $O(\log k)\cdot O(n^{2.805})$.

\end{enumerate}
\end{solution}

\item \label{triangle-detection} Let $G = (V,E)$ be an undirected graph. A \emph{triangle} in $G$ is a set of distinct vertices $u,v,w \in V$ that are pairwise connected by an edge.
Design a randomized algorithm that decides whether $G$ contains a triangle. You can assume that $G$ is represented by its \emph{adjacency matrix} $A \in \{0,1\}^{V \times V}$ where  $a_{uv} = 1$ if $uv \in E$ and otherwise $a_{uv} = 0$.  
It should satisfy the following.
\begin{enumerate}[i)] 
\item If $G$ does not contain a triangle, then it always asserts \emph{No}.
\item If the algorithm asserts \emph{Yes}, then $G$ contains a triangle.
\item The probability of the assertion \emph{No} provided that $G$ has a triangle should be bounded by $1/2$.
\item The algorithm should perform $O(n^{2.805})$ arithmetic operations on integers of size $O(\log \log n)$. 
\end{enumerate}

\begin{solution}
Let \(n := |V|\). First note that \(G\) contains a triangle if and only if there is an edge \(uv \in E\) such that $(A^2)_{uv} > 0$ since \((A^2)_{uv}\) counts the number of vertices \(w\) adjacent to both \(u\) and \(v\). 

\textbf{Algorithm.}
Let \(p_1,\dots,p_k\) be the first $k := \lceil 2 \log_2 n\rceil$ prime numbers. Do the following:

\begin{enumerate}[1).]
    \item Choose a prime \(p\) uniformly at random from \(\{p_1,\dots,p_k\}\).
    \item Compute the matrix \(B := A^2 \bmod p\) by using fast matrix multiplication over the field \(\mathbb{F}_p\).
    \item If there exists an edge \(uv \in E\) such that \(B_{uv} \neq 0\), then
    output \emph{Yes}. Otherwise output \emph{No}.
\end{enumerate}
\textbf{Correctness.}
\begin{enumerate}[i)]
  \item If \(G\) does not contain a triangle, i.e., for each edge \(uv \in E\), \(u\) and \(v\) have no common neighbor. Hence $B_{uv} = 0$ for every \(uv \in E\). Thus the algorithm always outputs \emph{No}.
  \item This statement is the same as i).
  \item Let \(u,v,w\) be a triangle in $G$ and consider the edge \(uv\). Then $1\leq B_{uv} \leq n$.
The algorithm fails to detect this edge if and only if
\[
B_{uv} \equiv 0 \pmod p.
\]
Since the number of distinct prime divisors of \(B_{uv}\) is at most \(\log_2 n\), one has
\[
\Pr[\text{Algorithm outputs \emph{No}} \mid G \text{ contains a triangle}\;uvw]
\leq \frac{\log_2 n}{\lceil2\log_2 n\rceil}
\leq \frac12.
\]
\item The \(k\)-th prime is \(O(k \log k)\) hence $p_k = O(\log n \log \log n)$.
Thus every arithmetic operation is performed on integers of size $O(\log p_k) = O(\log \log n)$.

Computing \(A^2 \bmod p\) can be done with \(O(n^{2.805})\) arithmetic operations over \(\mathbb{F}_p\) using fast matrix multiplication algorithms. 
Finally searching all the entries of $B$ can be done in \(O(n^2)\) time.
\end{enumerate}
\end{solution}

\item The \emph{node-edge incidence matrix} $A_G \in \{0, \pm 1\}^{V \times A}$ of a directed graph $G = (V,A)$   is defined as
  \begin{displaymath}
    a_{ua} = \begin{cases}
      1, & \text{if  } a= (x,u) \text{ for some } x \in V,\\
      -1, & \text{if  } a= (u,x) \text{ for some } x \in V,\\
0, & \text{otherwise.}
    \end{cases}    
  \end{displaymath}
  Show that $A_G$ is totally unimodular. 

\begin{solution}
We will prove by induction that $\det(B) \in \{0,\pm 1\}$ for any square submatrix $B$ of $A_G$. The base case is trivial. For the induction step, suppose that the claim holds for all square submatrices of size $\leq k-1$, and let $B$ be a square submatrix of size $k$.
If $B$ has a column with at most one nonzero entry, we are done since by expanding the determinant along that column gives $\det(B) \in \{0,\pm \det(B')\}$,
where $B'$ is a cofactor matrix of $B$ of size $k-1$. By induction hypothesis,
$\det(B') \in \{0,\pm 1\}$, hence $\det(B)\in \{0,\pm 1\}$.

Then it remains to show the case where every column of $B$ has exactly two nonzero entries. Since every column of the incidence matrix $A_G$ has at most
two nonzero entries, namely one $1$ and one $-1$, each column of $B$ must
contain both of these entries. Therefore, the sum of all rows of $B$ is zero. Hence the rows of $B$ are linearly dependent, i.e., $\det(B)=0$.
\end{solution}
\end{enumerate}


  

\end{document}

%%% Local Variables:
%%% mode: latex
%%% TeX-master: t
%%% End:
